\corrigesChapitre{15}{Mesurer la tension électrique}
\par\noindent\begin{minipage}{\linewidth}
\begin{corrige}{1}{La bonne grandeur}
\begin{questions}
\item La tension électrique.
\item Un voltmètre, souvent une fonction du multimètre.
\end{questions}
\end{corrige}
\end{minipage}
\par\noindent\begin{minipage}{\linewidth}
\begin{corrige}{2}{Deux connexions}
\begin{questions}
\item En dérivation aux deux bornes de la lampe.
\item Non : on relie les pointes à ses deux bornes selon le protocole.
\end{questions}
\pcFigureEx{figures/cor-voltmetre.pdf}{Les deux bornes du voltmètre sont reliées aux deux bornes de la lampe : on n’ouvre pas la branche principale.}
\end{corrige}
\end{minipage}
\par\noindent\begin{minipage}{\linewidth}
\begin{corrige}{3}{Un circuit ouvert}
\begin{questions}
\item Oui, entre ses bornes.
\item Non : sans trajet fermé extérieur, le courant permanent extérieur est nul dans le modèle.
\end{questions}
\pcFigureEx{figures/cor-pile-isolee.pdf}{Une tension peut exister sans courant dans un circuit extérieur fermé.}
\end{corrige}
\end{minipage}
\par\noindent\begin{minipage}{\linewidth}
\begin{corrige}{4}{Convertir}
\begin{questions}
\item $2\,500\unite{mV}$.
\item $0{,}600\unite V$.
\end{questions}
\end{corrige}
\end{minipage}
\par\noindent\begin{minipage}{\linewidth}
\begin{corrige}{5}{Une dérivation simple}
\begin{questions}
\item Chaque lampe a une tension de 4,5 V avec la même polarité.
\item Non : l’égalité des tensions en dérivation ne correspond pas à un partage par deux.
\end{questions}
\pcFigureEx{figures/cor-tensions-45.pdf}{Les lampes sont reliées aux mêmes deux nœuds : on ne divise pas la tension par le nombre de branches.}
\end{corrige}
\end{minipage}
\par\noindent\begin{minipage}{\linewidth}
\begin{corrige}{6}{Compléter une somme}
\begin{questions}
\item $U_2=9-4=5\unite V$.
\item $4+5=9\unite V$.
\end{questions}
\par\smallskip\noindent\textbf{Pour comprendre.}\enspace Dans la boucle étudiée, $U_G=U_1+U_2$. Donc $U_2=9-4=5\unite V$ ; on retrouve bien 9 V en additionnant les tensions des récepteurs.\par
\pcFigureEx{figures/cor-tensions-serie.pdf}{Dans cette boucle, les tensions des récepteurs se somment pour retrouver celle de la source.}
\end{corrige}
\end{minipage}
\par\noindent\begin{minipage}{\linewidth}
\begin{corrige}{7}{Trois récepteurs}
\begin{questions}
\item $U_G=1{,}2+2{,}0+2{,}8=6{,}0\unite V$.
\item Les récepteurs sont orientés de manière cohérente dans la boucle et les fils sont idéaux.
\end{questions}
\end{corrige}
\end{minipage}
\par\noindent\begin{minipage}{\linewidth}
\begin{corrige}{8}{Inverser les pointes}
\begin{questions}
\item $-3{,}0\unite V$.
\item Non : le signe traduit la polarité de la mesure.
\end{questions}
\pcFigureEx{figures/cor-voltmetre-inverse.pdf}{Inverser les pointes inverse le signe de la différence de potentiel mesurée, sans modifier sa valeur absolue.}
\end{corrige}
\end{minipage}
\par\noindent\begin{minipage}{\linewidth}
\begin{corrige}{9}{Un interrupteur idéal}
\begin{questions}
\item Une tension nulle.
\item Oui : une tension nulle sur un fil idéal n’impose pas un courant nul.
\end{questions}
\pcFigureEx{figures/cor-interrupteur-ferme.pdf}{Dans le modèle idéal, un interrupteur fermé a une tension nulle et peut laisser passer un courant.}
\end{corrige}
\end{minipage}
\par\noindent\begin{minipage}{\linewidth}
\begin{corrige}{10}{Même tension, autre courant}
\begin{questions}
\item Oui si elles sont reliées aux mêmes deux nœuds avec la même polarité.
\item Non : les intensités dépendent aussi des caractéristiques des lampes.
\end{questions}
\end{corrige}
\end{minipage}
\par\noindent\begin{minipage}{\linewidth}
\begin{corrige}{11}{Deux lampes non identiques}
\begin{questions}
\item Oui : 1,8 + 4,2 = 6 V.
\item Le partage égal n’est pas une loi générale ; les lampes n’ont pas forcément les mêmes caractéristiques.
\end{questions}
\par\smallskip\noindent\textbf{Pour comprendre.}\enspace En série, la loi impose la somme : $1{,}8+4{,}2=6{,}0\unite V$. Un partage égal ne peut être ajouté sans hypothèse supplémentaire sur les dipôles.\par
\pcFigureEx{figures/cor-tensions-18-42.pdf}{La somme vaut 6 V ; la loi ne demande pas que deux récepteurs différents se partagent la tension à parts égales.}
\end{corrige}
\end{minipage}
\par\noindent\begin{minipage}{\linewidth}
\begin{corrige}{12}{Une mesure en trop ?}
\begin{questions}
\item Non : les deux points appartiennent au même nœud idéal.
\item Non : pour la lampe, il faut relier les pointes à ses deux bornes distinctes.
\end{questions}
\pcFigureEx{figures/cor-voltmetre-meme-noeud.pdf}{Les deux pointes sont sur le même conducteur idéal ; elles ne sont pas aux deux bornes de la lampe.}
\end{corrige}
\end{minipage}
\par\noindent\begin{minipage}{\linewidth}
\begin{corrige}{13}{Un bilan mesuré}
\begin{questions}
\item $U_1+U_2=6{,}00\unite V$, soit un écart de $0{,}02\unite V$.
\item Comparer cet écart à la précision et à la stabilité des appareils ; les mesures peuvent être compatibles avec l’additivité.
\end{questions}
\par\smallskip\noindent\textbf{Pour comprendre.}\enspace $U_1+U_2=2{,}51+3{,}49=6{,}00\unite V$, soit un écart de 0,02 V à la source mesurée. Cet écart se discute avec l’appareil et les conditions, pas avec le seul nombre de décimales affiché.\par
\end{corrige}
\end{minipage}
\par\noindent\begin{minipage}{\linewidth}
\begin{corrige}{14}{Une confusion dangereuse}
\begin{questions}
\item Il risque un court-circuit par l’appareil de faible résistance.
\item Fonction tension continue, bornes V et COM, calibre adapté et schéma validé ; l’élève utilise seulement le matériel autorisé.
\end{questions}
\pcFigureEx{figures/cor-voltmetre.pdf}{Vérifier mode tension, bornes V/COM et calibre avant de relier le voltmètre en dérivation.}
\end{corrige}
\end{minipage}
\par\noindent\begin{minipage}{\linewidth}
\begin{corrige}{15}{Retrouver le montage}
\begin{questions}
\item Non : en série, leur somme ferait 12 V.
\item Une dérivation directe sur les deux bornes du générateur est compatible ; d’autres informations peuvent être nécessaires pour identifier tout le circuit.
\end{questions}
\pcFigureEx{figures/cor-serie-derivation-6V.pdf}{Deux tensions positives de 6 V ne peuvent pas s’ajouter pour donner une source de 6 V ; la dérivation directe est compatible.}
\end{corrige}
\end{minipage}
\par\noindent\begin{minipage}{\linewidth}
\phantomsection\label{cor-defi-15}\pcSolutionDefi{15}{Distinguer deux montages}\begin{questions}
\item A est compatible avec une série ; B avec une dérivation directe.
\item La série impose $I_1=I_2$, mais pas leur valeur numérique sans caractéristique supplémentaire.
\item Non : l’égalité des tensions ne suffit pas à imposer l’égalité des courants pour des dipôles différents.
\end{questions}
\pcFigureEx{figures/cor-serie-derivation-9V.pdf}{En A : 3 V + 6 V = 9 V ; en B, chaque récepteur reçoit les 9 V de la source.}
\end{minipage}
\par\noindent\begin{minipage}{\linewidth}
\phantomsection\label{cor-activite-15}\pcSousTitre{Pistes pour l’investigation}
En série, la somme des tensions des récepteurs est compatible avec la tension de la source ; en dérivation, les dipôles entre les mêmes nœuds ont la même tension orientée. Les bornes V et COM doivent correspondre aux polarités annoncées. Un signe négatif révèle une inversion de mesure, et non nécessairement un défaut du circuit.\par\end{minipage}
